\documentclass[10pt,a4paper]{amsart}
\usepackage{amsmath,amssymb,amsthm,graphicx}
% Retained preamble for the Stacks Project chapter {CH} (shared project preamble).
% The chapter file pinned-original-source.tex only carries \input{{preamble}} on its first line,
% so no preamble text lives inside the chapter itself.  The Stacks Project uses one single shared
% preamble for every chapter: the snapshot stages it three times (sources/stacks-chow/preamble.tex,
% sources/stacks-etale-cohomology/preamble.tex and focused-seeds/AGP/sources/stacks-snapshot-20261011/
% preamble.tex) and those copies are identical in size.  The block below is copied mechanically,
% line for line, from sources/stacks-chow/preamble.tex lines 14-18 and 154-217: the Xy-pic setup and
% the native theorem declarations together with the mathematical macros.  Nothing was added, removed
% or redefined; the wrapper class and the amsmath/amssymb/amsthm/graphicx base packages are supplied
% by the assembler.
% For commutative diagrams we use Xy-pic
% Chapter: stacks-spaces-more-morphisms.tex -- More on Morphisms of Spaces.
\usepackage[all]{xy}

% We use 2cell for 2-commutative diagrams.
\xyoption{2cell}

% Theorem environments.
%
\theoremstyle{plain}
\newtheorem{theorem}[subsection]{Theorem}
\newtheorem{proposition}[subsection]{Proposition}
\newtheorem{lemma}[subsection]{Lemma}

\theoremstyle{definition}
\newtheorem{definition}[subsection]{Definition}
\newtheorem{example}[subsection]{Example}
\newtheorem{exercise}[subsection]{Exercise}
\newtheorem{situation}[subsection]{Situation}

\theoremstyle{remark}
\newtheorem{remark}[subsection]{Remark}
\newtheorem{remarks}[subsection]{Remarks}

\numberwithin{equation}{subsection}

% Macros
%
\def\lim{\mathop{\mathrm{lim}}\nolimits}
\def\colim{\mathop{\mathrm{colim}}\nolimits}
\def\Spec{\mathop{\mathrm{Spec}}}
\def\Hom{\mathop{\mathrm{Hom}}\nolimits}
\def\Ext{\mathop{\mathrm{Ext}}\nolimits}
\def\SheafHom{\mathop{\mathcal{H}\!\mathit{om}}\nolimits}
\def\SheafExt{\mathop{\mathcal{E}\!\mathit{xt}}\nolimits}
\def\Sch{\mathit{Sch}}
\def\Mor{\mathop{\mathrm{Mor}}\nolimits}
\def\Ob{\mathop{\mathrm{Ob}}\nolimits}
\def\Sh{\mathop{\mathit{Sh}}\nolimits}
\def\NL{\mathop{N\!L}\nolimits}
\def\CH{\mathop{\mathrm{CH}}\nolimits}
\def\proetale{{pro\text{-}\acute{e}tale}}
\def\etale{{\acute{e}tale}}
\def\QCoh{\mathit{QCoh}}
\def\Ker{\mathop{\mathrm{Ker}}}
\def\Im{\mathop{\mathrm{Im}}}
\def\Coker{\mathop{\mathrm{Coker}}}
\def\Coim{\mathop{\mathrm{Coim}}}

% Boxtimes
%
\DeclareMathSymbol{\boxtimes}{\mathbin}{AMSa}{"02}

%
% Macros for moduli stacks/spaces
%
\def\QCohstack{\mathcal{QC}\!\mathit{oh}}
\def\Cohstack{\mathcal{C}\!\mathit{oh}}
\def\Spacesstack{\mathcal{S}\!\mathit{paces}}
\def\Quotfunctor{\mathrm{Quot}}
\def\Hilbfunctor{\mathrm{Hilb}}
\def\Curvesstack{\mathcal{C}\!\mathit{urves}}
\def\Polarizedstack{\mathcal{P}\!\mathit{olarized}}
\def\Complexesstack{\mathcal{C}\!\mathit{omplexes}}
% \Pic is the operator that assigns to X its picard group, usage \Pic(X)
% \Picardstack_{X/B} denotes the Picard stack of X over B
% \Picardfunctor_{X/B} denotes the Picard functor of X over B
\def\Pic{\mathop{\mathrm{Pic}}\nolimits}
\def\Picardstack{\mathcal{P}\!\mathit{ic}}
\def\Picardfunctor{\mathrm{Pic}}
\def\Deformationcategory{\mathcal{D}\!\mathit{ef}}
\usepackage{cleveref}
\crefname{originalreference}{original reference}{original references}
\Crefname{originalreference}{Original reference}{Original references}
\makeatletter
\newlabel{decent-spaces-lemma-generalizations-lift-flat}{{0B7W}{source}{}{sourcepointer-1}{}}
\newlabel{decent-spaces-lemma-generalizations-lift-flat@cref}{{[originalreference][0][]0B7W}{[1][0][]source}{}{}{}}
\makeatother
\crefname{originalreference}{original reference}{original references}
\Crefname{originalreference}{Original reference}{Original references}
\makeatletter
\newlabel{decent-spaces-lemma-specialization}{{03K2}{source}{}{sourcepointer-2}{}}
\newlabel{decent-spaces-lemma-specialization@cref}{{[originalreference][0][]03K2}{[1][0][]source}{}{}{}}
\makeatother
\crefname{originalreference}{original reference}{original references}
\Crefname{originalreference}{Original reference}{Original references}
\makeatletter
\newlabel{more-algebra-lemma-lift-from-fibre-relatively-perfect}{{0DJH}{source}{}{sourcepointer-3}{}}
\newlabel{more-algebra-lemma-lift-from-fibre-relatively-perfect@cref}{{[originalreference][0][]0DJH}{[1][0][]source}{}{}{}}
\makeatother
\crefname{originalreference}{original reference}{original references}
\Crefname{originalreference}{Original reference}{Original references}
\makeatletter
\newlabel{more-algebra-lemma-no-change-tor-amplitude}{{0DJF}{source}{}{sourcepointer-4}{}}
\newlabel{more-algebra-lemma-no-change-tor-amplitude@cref}{{[originalreference][0][]0DJF}{[1][0][]source}{}{}{}}
\makeatother
\crefname{originalreference}{original reference}{original references}
\Crefname{originalreference}{Original reference}{Original references}
\makeatletter
\newlabel{perfect-lemma-affine-compare-bounded}{{06Z0}{source}{}{sourcepointer-5}{}}
\newlabel{perfect-lemma-affine-compare-bounded@cref}{{[originalreference][0][]06Z0}{[1][0][]source}{}{}{}}
\makeatother
\crefname{originalreference}{original reference}{original references}
\Crefname{originalreference}{Original reference}{Original references}
\makeatletter
\newlabel{perfect-lemma-pseudo-coherent-affine}{{08E7}{source}{}{sourcepointer-6}{}}
\newlabel{perfect-lemma-pseudo-coherent-affine@cref}{{[originalreference][0][]08E7}{[1][0][]source}{}{}{}}
\makeatother
\crefname{originalreference}{original reference}{original references}
\Crefname{originalreference}{Original reference}{Original references}
\makeatletter
\newlabel{sites-cohomology-lemma-tor-amplitude-stalk}{{0DJJ}{source}{}{sourcepointer-7}{}}
\newlabel{sites-cohomology-lemma-tor-amplitude-stalk@cref}{{[originalreference][0][]0DJJ}{[1][0][]source}{}{}{}}
\makeatother
\crefname{originalreference}{original reference}{original references}
\Crefname{originalreference}{Original reference}{Original references}
\makeatletter
\newlabel{spaces-perfect-lemma-derived-quasi-coherent-small-etale-site}{{071Q}{source}{}{sourcepointer-8}{}}
\newlabel{spaces-perfect-lemma-derived-quasi-coherent-small-etale-site@cref}{{[originalreference][0][]071Q}{[1][0][]source}{}{}{}}
\makeatother
\crefname{originalreference}{original reference}{original references}
\Crefname{originalreference}{Original reference}{Original references}
\makeatletter
\newlabel{spaces-perfect-lemma-descend-pseudo-coherent}{{08HE}{source}{}{sourcepointer-9}{}}
\newlabel{spaces-perfect-lemma-descend-pseudo-coherent@cref}{{[originalreference][0][]08HE}{[1][0][]source}{}{}{}}
\makeatother
\crefname{originalreference}{original reference}{original references}
\Crefname{originalreference}{Original reference}{Original references}
\makeatletter
\newlabel{spaces-perfect-lemma-pseudo-coherent}{{08JL}{source}{}{sourcepointer-10}{}}
\newlabel{spaces-perfect-lemma-pseudo-coherent@cref}{{[originalreference][0][]08JL}{[1][0][]source}{}{}{}}
\makeatother
\crefname{originalreference}{original reference}{original references}
\Crefname{originalreference}{Original reference}{Original references}
\makeatletter
\newlabel{spaces-properties-theorem-exactness-stalks}{{04K5}{source}{}{sourcepointer-11}{}}
\newlabel{spaces-properties-theorem-exactness-stalks@cref}{{[originalreference][0][]04K5}{[1][0][]source}{}{}{}}
\makeatother
\title{Relative perfectness is local for the etale topology on henselian pairs}
\author{The Stacks Project Authors}
\date{stacks-project:spaces-more-morphisms.tex:local-snapshot-20261011; GNU FDL 1.2 or later}
\begin{document}
\maketitle
\noindent\textit{Editorial note.} Counted span is the author's own proof environment, source lines 13181-13261, immediately following the statement at 13170-13179.\par\medskip
\noindent\textit{Editorial note.} Same-name pair note: this is the more general (algebraic spaces) side; the schemes side is not counted beside it.\par\medskip
\noindent\textit{Editorial note.} Same-chapter pointers (characterisation, thickenings) are recorded in assembly.external\_references.\par\medskip
\noindent\textit{Editorial note.} Steps the author leaves to the reader are flagged here for the reviewer: source line(s) 13234 carry an "omitted"/"follows easily"/"immediate from the definitions" phrase. Those sentences are retained verbatim; the displayed argument around them stays inside the counted span.\par\medskip
\noindent\textit{Allowed standard prerequisites (editorial).}\begin{itemize}
\item Characterisation of relatively perfect morphisms. chapter lemma lemma-characterize-relatively-perfect, cited in the proof. Supplies the criterion that is checked etale-locally.
\item Thickenings and pseudo-coherence. chapter lemmas lemma-thickening-pseudo-coherent and lemma-thickening-relatively-perfect, cited in the proof. Control the behaviour of the property along the thickenings appearing in the henselian setup.
\item Henselian pairs and etale localization. the notion of a henselian pair used in the statement together with the etale-local reduction in the proof. Reduces the verification to the affine henselian case where the criterion applies.
\end{itemize}\medskip
\ifdefined\hypertarget\hypertarget{sourcepointer-1}{}\fi
\noindent\textit{Source reference.} Original reference 0B7W: focused-seeds/geometry/S9-spaces-henselian-relatively-perfect-spaces/ref-tags.txt, source lines 5--5. Navigation pointer only: the counted proof cites this statement by \textbackslash{}ref as a standing input; the deduction the counted proof has to make is carried out inside the counted span, and no part of that deduction is delegated to this pointer.\par\medskip
\ifdefined\hypertarget\hypertarget{sourcepointer-2}{}\fi
\noindent\textit{Source reference.} Original reference 03K2: focused-seeds/geometry/S9-spaces-henselian-relatively-perfect-spaces/ref-tags.txt, source lines 6--6. Navigation pointer only: the counted proof cites this statement by \textbackslash{}ref as a standing input; the deduction the counted proof has to make is carried out inside the counted span, and no part of that deduction is delegated to this pointer.\par\medskip
\ifdefined\hypertarget\hypertarget{sourcepointer-3}{}\fi
\noindent\textit{Source reference.} Original reference 0DJH: focused-seeds/geometry/S9-spaces-henselian-relatively-perfect-spaces/ref-tags.txt, source lines 7--7. Navigation pointer only: the counted proof cites this statement by \textbackslash{}ref as a standing input; the deduction the counted proof has to make is carried out inside the counted span, and no part of that deduction is delegated to this pointer.\par\medskip
\ifdefined\hypertarget\hypertarget{sourcepointer-4}{}\fi
\noindent\textit{Source reference.} Original reference 0DJF: focused-seeds/geometry/S9-spaces-henselian-relatively-perfect-spaces/ref-tags.txt, source lines 8--8. Navigation pointer only: the counted proof cites this statement by \textbackslash{}ref as a standing input; the deduction the counted proof has to make is carried out inside the counted span, and no part of that deduction is delegated to this pointer.\par\medskip
\ifdefined\hypertarget\hypertarget{sourcepointer-5}{}\fi
\noindent\textit{Source reference.} Original reference 06Z0: focused-seeds/geometry/S9-spaces-henselian-relatively-perfect-spaces/ref-tags.txt, source lines 9--9. Navigation pointer only: the counted proof cites this statement by \textbackslash{}ref as a standing input; the deduction the counted proof has to make is carried out inside the counted span, and no part of that deduction is delegated to this pointer.\par\medskip
\ifdefined\hypertarget\hypertarget{sourcepointer-6}{}\fi
\noindent\textit{Source reference.} Original reference 08E7: focused-seeds/geometry/S9-spaces-henselian-relatively-perfect-spaces/ref-tags.txt, source lines 10--10. Navigation pointer only: the counted proof cites this statement by \textbackslash{}ref as a standing input; the deduction the counted proof has to make is carried out inside the counted span, and no part of that deduction is delegated to this pointer.\par\medskip
\ifdefined\hypertarget\hypertarget{sourcepointer-7}{}\fi
\noindent\textit{Source reference.} Original reference 0DJJ: focused-seeds/geometry/S9-spaces-henselian-relatively-perfect-spaces/ref-tags.txt, source lines 11--11. Navigation pointer only: the counted proof cites this statement by \textbackslash{}ref as a standing input; the deduction the counted proof has to make is carried out inside the counted span, and no part of that deduction is delegated to this pointer.\par\medskip
\ifdefined\hypertarget\hypertarget{sourcepointer-8}{}\fi
\noindent\textit{Source reference.} Original reference 071Q: focused-seeds/geometry/S9-spaces-henselian-relatively-perfect-spaces/ref-tags.txt, source lines 12--12. Navigation pointer only: the counted proof cites this statement by \textbackslash{}ref as a standing input; the deduction the counted proof has to make is carried out inside the counted span, and no part of that deduction is delegated to this pointer.\par\medskip
\ifdefined\hypertarget\hypertarget{sourcepointer-9}{}\fi
\noindent\textit{Source reference.} Original reference 08HE: focused-seeds/geometry/S9-spaces-henselian-relatively-perfect-spaces/ref-tags.txt, source lines 13--13. Navigation pointer only: the counted proof cites this statement by \textbackslash{}ref as a standing input; the deduction the counted proof has to make is carried out inside the counted span, and no part of that deduction is delegated to this pointer.\par\medskip
\ifdefined\hypertarget\hypertarget{sourcepointer-10}{}\fi
\noindent\textit{Source reference.} Original reference 08JL: focused-seeds/geometry/S9-spaces-henselian-relatively-perfect-spaces/ref-tags.txt, source lines 14--14. Navigation pointer only: the counted proof cites this statement by \textbackslash{}ref as a standing input; the deduction the counted proof has to make is carried out inside the counted span, and no part of that deduction is delegated to this pointer.\par\medskip
\ifdefined\hypertarget\hypertarget{sourcepointer-11}{}\fi
\noindent\textit{Source reference.} Original reference 04K5: focused-seeds/geometry/S9-spaces-henselian-relatively-perfect-spaces/ref-tags.txt, source lines 15--15. Navigation pointer only: the counted proof cites this statement by \textbackslash{}ref as a standing input; the deduction the counted proof has to make is carried out inside the counted span, and no part of that deduction is delegated to this pointer.\par\medskip
\begin{lemma}
\label{lemma-henselian-relatively-perfect}
Let $(R, I)$ be a pair consisting of a ring and an ideal $I$
contained in the Jacobson radical. Set $S = \Spec(R)$ and $S_0 = \Spec(R/I)$.
Let $X$ be an algebraic space over $R$ whose structure morphism
$f : X \to S$ is proper, flat, and of finite presentation.
Denote $X_0 = S_0 \times_S X$. Let $E \in D(\mathcal{O}_X)$
be pseudo-coherent. If the derived restriction $E_0$ of $E$
to $X_0$ is $S_0$-perfect, then $E$ is $S$-perfect.
\end{lemma}

\begin{proof}
Choose a surjective \'etale morphism $U \to X$ with $U$ affine.
Choose a closed immersion $U \to \mathbf{A}^d_S$.
Set $U_0 = S_0 \times_S U$.
The complex $E_0|_{U_0}$ has tor amplitude
in $[a, b]$ for some $a, b \in \mathbf{Z}$.
Let $\overline{x}$ be a geometric point of $X$.
We will show that the tor amplitude of
$E_{\overline{x}}$ over $R$ is in $[a - d, b]$.
This will finish the proof as the tor amplitude can be
read off from the stalks by
Cohomology on Sites, Lemma \ref{sites-cohomology-lemma-tor-amplitude-stalk}
and Properties of Spaces, Theorem
\ref{spaces-properties-theorem-exactness-stalks}.

\medskip\noindent
Let $x \in |X|$ be the point determined by $\overline{x}$.
Recall that $|X| \to |S|$ is closed (by definition of proper morphisms).
Since $I$ is contained in the Jacobson radical, any nonempty closed
subset of $S$ contains a point of the closed subscheme $S_0$.
Hence we can find a specialization $x \leadsto x_0$ in $|X|$
with $x_0 \in |X_0|$. Choose $u_0 \in U_0$ mapping to $x_0$.
By Decent Spaces, Lemma \ref{decent-spaces-lemma-generalizations-lift-flat}
(or by Decent Spaces, Lemma \ref{decent-spaces-lemma-specialization}
which applies directly to \'etale morphisms)
we find a specialization $u \leadsto u_0$ in $U$
such that $u$ maps to $x$. We may lift $\overline{x}$
to a geometric point $\overline{u}$ of $U$ lying over $u$.
Then we have $E_{\overline{x}} = (E|_U)_{\overline{u}}$.

\medskip\noindent
Write $U = \Spec(A)$. Then $A$ is a flat, finitely presented
$R$-algebra which is a quotient of a polynomial $R$-algebra in
$d$-variables. The restriction $E|_U$ corresponds
(by Derived Categories of Spaces, Lemmas
\ref{spaces-perfect-lemma-pseudo-coherent},
\ref{spaces-perfect-lemma-derived-quasi-coherent-small-etale-site}, and
\ref{spaces-perfect-lemma-descend-pseudo-coherent}
and
Derived Categories of Schemes, Lemma
\ref{perfect-lemma-affine-compare-bounded} and
\ref{perfect-lemma-pseudo-coherent-affine})
to a pseudo-coherent object $K$ of $D(A)$.
Observe that $E_0$ corresponds to $K \otimes_A^\mathbf{L} A/IA$.
Let $\mathfrak q \subset \mathfrak q_0 \subset A$ be the prime
ideals corresponding to $u \leadsto u_0$.
Then
$$
E_{\overline{x}} =
(E|_U)_{\overline{u}} =
E_u \otimes_{\mathcal{O}_{U, u}}^\mathbf{L} \mathcal{O}_{U, \overline{u}} =
K_{\mathfrak q} \otimes_{A_\mathfrak q}^\mathbf{L} A_{\mathfrak q}^{sh}
$$
(some details omitted). Since $A_\mathfrak q \to A_\mathfrak q^{sh}$
is flat, the tor amplitude of this as an $R$-module is the same as
the tor amplitude of $K_\mathfrak q$ as an $R$-module
(More on Algebra, Lemma \ref{more-algebra-lemma-no-change-tor-amplitude}).
Also, $K_{\mathfrak q}$ is a localization of $K_{\mathfrak q_0}$.
Hence it suffices to show that $K_{\mathfrak q_0}$ has tor amplitude in
$[a - d, b]$ as a complex of $R$-modules.

\medskip\noindent
Let $I \subset \mathfrak p_0 \subset R$ be the prime
ideal corresponding to $f(x_0)$. Then we have
\begin{align*}
K \otimes_R^\mathbf{L} \kappa(\mathfrak p_0)
& =
(K \otimes_R^\mathbf{L} R/I) \otimes_{R/I}^\mathbf{L}
\kappa(\mathfrak p_0) \\
& =
(K \otimes_A^\mathbf{L} A/IA) \otimes_{R/I}^\mathbf{L} \kappa(\mathfrak p_0)
\end{align*}
the second equality because $R \to A$ is flat.
By our choice of $a, b$ this complex has cohomology
only in degrees in the interval $[a, b]$.
Thus we may finally apply
More on Algebra, Lemma
\ref{more-algebra-lemma-lift-from-fibre-relatively-perfect}
to $R \to A$, $\mathfrak q_0$, $\mathfrak p_0$ and $K$
to conclude.
\end{proof}

\end{document}
